% -----------------------------------------------------------------------------
% -----------------------------------------------------------------------------

\begin{exbox}[200\,keV 电子穿过 800\,nm 高斯近场]
\eqnrefs{eq:pinem_standard_beta_pulsed} 已经把经典 PINEM 的全部场依赖压缩成沿电子轨迹的同步 Fourier 投影。这个积分衡量的是电子轨迹对同步 Fourier 分量的选择。取一个最简单但非平凡的纵向近场包络
\begin{equation*}
 E_{\omega,z}(z)=E_{\omega,0}\exp\!\left(-\frac{z^2}{2L^2}\right),
 \qquad f(t)=1,
\end{equation*}
其中 $E_{\omega,0}$ 是正频电场复振幅的峰值（实场余弦峰值为其模长的两倍），$L$ 是近场沿轨迹的均方根空间宽度。电子初始动能取
\begin{equation*}
 K_0=200\,\mathrm{keV},
\end{equation*}
光学波长、近场峰值和空间宽度分别取
\begin{equation*}
 \lambda=800\,\mathrm{nm},\qquad
 E_{\omega,0}=20\,\mathrm{MV\,m^{-1}},\qquad
 L=100\,\mathrm{nm}.
\end{equation*}

先由自由电子相对论运动学得到
\begin{align*}
 \gamma_0
 &=1+\frac{K_0}{\me c^2}
 =1.39139,\\
 \frac{v_0}{c}
 &=\sqrt{1-\gamma_0^{-2}}
 =0.69531,\\
 v_0&=2.0845\times10^8\,\mathrm{m\,s^{-1}}.
\end{align*}
光学角频率和单光子能量则为
\begin{align*}
 \omega&=\frac{2\pi c}{\lambda}
 =2.3546\times10^{15}\,\mathrm{s^{-1}},\\
 \hbar\omega&=1.5498\,\mathrm{eV}.
\end{align*}
在这个例子中，\eqnrefs{eq:pinem_standard_beta_pulsed} 变成一个标准 Gaussian Fourier 积分：
\begin{align}
 \beta
 &=\frac{\qe E_{\omega,0}}{\hbar\omega}
 \int_{-\infty}^{\infty}\dd z\,
 \exp\!\left(-\frac{z^2}{2L^2}\right)
 \exp\!\left(-\ii\frac{\omega z}{v_0}\right)\notag\\
 &=\frac{\qe E_{\omega,0}\sqrt{2\pi}L}{\hbar\omega}
 \exp\!\left[-\frac12\left(\frac{\omega L}{v_0}\right)^2\right].
 \label{eq:pinem-gaussian-beta-dp65}
\end{align}
这里
\begin{equation*}
 \varphi_{\rm pm}\equiv\frac{\omega L}{v_0}=1.1296
\end{equation*}
是电子飞越近场包络时累计的光学相位尺度，并由\eqnrefs{eq:pinem_standard_beta_pulsed} 中 Fourier 相位的无量纲组合直接确定。数值代入后
\begin{equation}
 \boxed{|\beta|=1.7092,\qquad 2|\beta|=3.4183.}
 \label{eq:pinem-gaussian-beta-number-dp65}
\end{equation}

于是理想低反冲边带概率不需要再解微分方程，直接是
\begin{equation*}
 P_\ell=J_\ell^2(3.4183).
\end{equation*}
几个主要边带为
\begin{center}
\begin{tabular}{c|cccccc}
 $|\ell|$ & 0 & 1 & 2 & 3 & 4 & 5\\\hline
 $P_\ell$ & 0.1351 & 0.02944 & 0.2189 & 0.1413 & 0.03685 & 0.00538
\end{tabular}
\end{center}
这里表中 $|\ell|\ge1$ 的数值是\emph{单侧}边带概率；因为当前纯相位调制满足 $P_{+\ell}=P_{-\ell}$，总概率仍按两侧共同归一化。输出能量峰位为
\begin{equation*}
 K_\ell\simeq K_0+\ell\hbar\omega
 =200\,\mathrm{keV}+1.5498\ell\,\mathrm{eV}.
\end{equation*}
Bessel 权重并不随 $|\ell|$ 单调衰减；例如这里 $|\ell|=2$ 比 $|\ell|=1$ 更强。这不是异常，而是不同量子交换路径在同一整数能量阶上的相干衍射结果。

利用恒等式
\begin{equation*}
 \sum_{\ell=-\infty}^{\infty}\ell^2J_\ell^2(x)=\frac{x^2}{2},
\end{equation*}
还可以不逐项求和就得到电子边带阶数的均方根宽度
\begin{align*}
 \sqrt{\langle\ell^2\rangle}
 &=\sqrt{2}\,|\beta|=2.417,\\
 \sigma_E
 &=\hbar\omega\sqrt{\langle\ell^2\rangle}
 =3.75\,\mathrm{eV}.
\end{align*}
这给出了一个很有用的尺度判断：$|\beta|$ 不是“平均吸收光子数”，而是控制相干能量衍射宽度的脉冲面积。

\begin{center}
\includegraphics[width=0.78\textwidth]{pinem_gaussian_example_dp65.png}\\[-0.2em]
{\small 计算图：200\,keV、800\,nm 高斯近场的离散边带概率。每一根竖线相隔 $\hbar\omega=1.5498\,\mathrm{eV}$，全部权重由同一个 $|\beta|=1.7092$ 决定。}
\end{center}

最后必须检查这组参数是否真的允许使用 Bessel 极限。局域最小耦合小参数为
\begin{equation*}
 \frac{|\qe|E_{\omega,0}/\omega}{p_0}
 =5.15\times10^{-6},
 \qquad p_0=\gamma_0\me v_0,
\end{equation*}
说明单点场强对载波动量只是微扰；但沿整个作用区相干积分后 $|\beta|$ 已经达到 $O(1)$。这两个量衡量的是不同问题，不能用“局域相互作用很弱”推断“最终只会出现一阶边带”。

再检查反冲。同步纵向波数为
\begin{equation*}
 \kappa_z=\frac{\omega}{v_0}=1.1296\times10^7\,\mathrm{m^{-1}},
\end{equation*}
纵向相对论有效质量为 $m_\parallel=\gamma_0^3\me=2.6937\me$，所以
\begin{align*}
 \omega_{\rm r}
 &=\frac{\hbar\kappa_z^2}{2m_\parallel}
 =2.742\times10^9\,\mathrm{s^{-1}},\\
 \hbar\omega_{\rm r}
 &=1.80\,\mu\mathrm{eV}.
\end{align*}
若用 $T_{\rm int}=2L/v_0=0.959\,\mathrm{fs}$ 作为作用时间尺度，并取实际显著占据窗口 $|\ell|\le5$，则\eqnrefs{eq:pinem_recoil_phase_condition} 给
\begin{equation}
 \boxed{
 \epsilon_{\rm rec}^{(T)}
 =(2\times5+1)\omega_{\rm r}T_{\rm int}
 =2.89\times10^{-5}.}
 \label{eq:pinem-gaussian-recoil-number-dp65}
\end{equation}
因此这组参数处在非常干净的低反冲区。反过来，在相同电子能量和光学频率下，如果周期结构把相互作用延长到 $T=1\,\mathrm{ps}$，同一参数已增至约 $3.0\times10^{-2}$；到 $10\,\mathrm{ps}$ 则约为 $0.30$。所以反冲是否重要取决于\emph{累计相位可分辨性}，而不只取决于单次交换的 $\hbar\omega/E_0$。
\end{exbox}
